\begin{abstract}
Language models deployed as agents produce a continual stream of interactions
whose outcomes can often be verified automatically, yet this experience is
typically discarded: weights are frozen at deployment, and the model repeats
its mistakes. Existing approaches to learning from deployment experience
either accumulate it in external memory, leaving the model itself unchanged,
or optimize the sparse outcome reward with reinforcement learning, which is
difficult to stabilize on long, heterogeneous interaction streams. We propose
\method{} (online self-distillation fine-tuning), which converts each verified
success into a weight update as it occurs. Concretely, verified outputs are
stored in a retrieval memory; the frozen base model, conditioned on a
problem's own verified answer as a hint, then serves as a teacher for the
deployed model, and its hint-conditioned token distribution is distilled into
LoRA adapters with a forward KL objective. The method requires no ground-truth
labels, no stronger teacher model, and no feedback beyond binary verification.
\method{} attains the best score on all seven StreamBench tasks, outperforming
retrieval-based in-context learning and REINFORCE-style baselines. On a fused
stream of 4{,}219 problems spanning three domains, a single continually
trained model outperforms three per-domain specialist models ($+3.4$ points on
code generation), whereas REINFORCE++ collapses under the induced distribution
shift. Moreover, at matched data and oracle budget, single-pass online updates
outperform multi-epoch offline training on the same experience. These results
suggest that verified deployment experience can serve directly as supervision
for continual model improvement.
\end{abstract}
